\section{The Grounded-Feedback Framework}
\label{sec:framework}

Our contribution is not a new self-correction method but an account of \emph{why} interaction is a separate axis, \emph{when} it helps (a feedback-actionability taxonomy), and a previously-missing requirement: that the \emph{evaluation} be grounded too, without which the effect is unmeasurable.

\subsection{A feedback taxonomy by grounding}
We classify the reviewer's feedback channel by how much external, task-conditional information it injects:
\begin{itemize}[leftmargin=1.4em,itemsep=1pt,topsep=2pt]
\item \textbf{Type 0 -- none.} Single-shot, no feedback.
\item \textbf{Type 1 -- ungrounded critique.} An LLM/VLM critiques the artifact from its weights (``looks wrong''). No execution, no measurement.
\item \textbf{Type 2 -- static analysis.} Linters/type-checkers: signal about the artifact's \emph{form}, not its behavior.
\item \textbf{Type 3 -- grounded.} The artifact is exercised and the result observed: \textbf{3a} execution (run tests, read tracebacks), \textbf{3b} rendered geometry (measure the DOM), \textbf{3c} temporal (sample animation frames), \textbf{3d} factual (search and verify claims).
\end{itemize}

\subsection{Why grounding breaks the ceiling: a DPI bound}
Let $A^\star$ be the (task-determined) correct artifact and $A$ a candidate. The proposer is a channel from its inputs $\Theta$ (weights) and prompt $X$ to $A$. A \emph{reviewer} that only reads $A$ and reasons forms the Markov chain $A^\star \to (\Theta,X) \to A \to R$, where $R$ is its critique. By the data-processing inequality,
\[
I(A^\star; R)\ \le\ I(A^\star; (\Theta,X)),
\]
so \emph{no} amount of reasoning or ungrounded review (Type 0--2) can recover more information about $A^\star$ than is already in $(\Theta,X)$. This bounds reasoning-only scaling. Sampling is bounded differently but no less tightly: for any finite $N$, best-of-$N$ can only \emph{select} among the $N$ candidates the proposer actually draws, so even an oracle (grounded) verifier cannot return a correct artifact the proposer is too unlikely to sample within the budget. Its practical ceiling is the high-probability support of the proposer's own output distribution, not a new information channel. Both ceilings appear empirically once we measure them (\Cref{fig:dpi} in \Cref{sec:interaction}): on hard code, both reasoning-only and oracle best-of-$N$ saturate well short of completion. The inequality bounds the \emph{information channel}, not achievable quality directly; we use it to \emph{motivate} the saturation we then measure and treat as the evidence.

Grounded feedback (Type 3) escapes both ceilings by feeding an environment observation $E=g(A,\,\mathrm{world})$ (a test outcome, a measured bounding box, a search verdict) back into \emph{generation}. Conditioning the next proposal on $E$ lets the proposer synthesize a candidate it would not otherwise have sampled: formally $I(A^\star;(R,E))$ can exceed $I(A^\star;(\Theta,X))$ because $E$ carries information obtained by \emph{running} the artifact against reality. This is the crucial difference from selection (best-of-$N$ cannot create a missing candidate) and from reasoning (which adds no information), and it is why interaction is a distinct axis rather than a constant-factor speedup.

\paragraph{The symmetric claim: ungrounded evaluation cannot measure the gain.} The same inequality applies to the \emph{scorer}. A metric $M$ that is itself a function of $(\Theta,X)$ (e.g.\ a VLM judging a screenshot, whose view of $A$ is a lossy, cropped image) satisfies $I(A^\star; M)\le I(A^\star;(\Theta,X))$ and cannot certify an improvement that lives in the part of $A$ it cannot observe. Grounding the evaluation (measure the artifact directly) is therefore not optional polish: it is a precondition for \emph{detecting} interaction scaling on any modality whose quality is not fully visible to the default judge. \Cref{sec:grounded_eval} shows this is exactly the situation for rendered visual artifacts.

\subsection{Predictions}
The framework makes three testable predictions, each tested below. \textbf{(P1)} Reasoning- and sampling-only scaling saturate strictly below the interaction ceiling (\Cref{sec:interaction}). \textbf{(P2)} Within interaction, more actionable grounded channels lift more, and the \emph{execution} of the artifact, not the extra review pass, is the active ingredient (\Cref{sec:interaction}). \textbf{(P3)} On modalities whose quality is invisible to the default judge, the gain is only detectable with a grounded metric, and an ungrounded reviewer can even regress quality (\Cref{sec:grounded_eval}).
